\documentclass[11pt,a4paper]{article}
\usepackage[T1]{fontenc}
\usepackage{lmodern}
\usepackage[margin=27mm]{geometry}
\usepackage{amsmath,amssymb,amsthm,mathtools}
\usepackage{microtype}
\usepackage{booktabs,array}
\usepackage{enumitem}
\usepackage[hidelinks]{hyperref}
\setlength{\parskip}{5pt plus 1pt}
\setlength{\parindent}{0pt}
\setlength{\emergencystretch}{3em}
\setlist{itemsep=2pt,topsep=3pt}
\raggedbottom
\newtheorem{theorem}{Theorem}[section]
\newtheorem{lemma}[theorem]{Lemma}
\newtheorem{corollary}[theorem]{Corollary}
\newtheorem{proposition}[theorem]{Proposition}
\theoremstyle{definition}
\newtheorem{definition}[theorem]{Definition}
\newtheorem{example}[theorem]{Example}
\theoremstyle{remark}
\newtheorem{remark}[theorem]{Remark}
\newcommand{\Z}{\mathbb{Z}}
\newcommand{\N}{\mathbb{N}}
\newcommand{\Q}{\mathbb{Q}}
\newcommand{\F}{\mathbb{F}}
\newcommand{\ps}[1]{#1[[x]]}
\newcommand{\coeff}[2]{[x^{#1}]\,#2}
\newcommand{\hl}[1]{\textbf{#1}}
\newcommand{\para}[1]{\par\medskip\noindent\textbf{#1}\enspace}
\hypersetup{pdftitle={Congruences modulo 4 for the series of OEIS A389542},pdfauthor={Dev D. Goyal},pdfsubject={OEIS A389542}}
\title{Congruences modulo $4$ for the series with $A\bigl(A(x)^2-x^2\bigr)=4A(x)^3$\\[6pt]\large A proof of a conjecture in OEIS A389542}
\author{Dev D. Goyal}
\date{October 2026}
\begin{document}
\maketitle
\begin{abstract}
The sequence A389542 of the On-Line Encyclopedia of Integer Sequences (OEIS) lists the coefficients of the power series
$A(x)=x+2x^2+10x^3+64x^4+\cdots$ that satisfies $A\bigl(A(x)^2-x^2\bigr)=4A(x)^3$. Paul~D.~Hanna conjectured that, for
$n>1$, the coefficient $a(n)$ leaves remainder $2$ on division by $4$ when $n-1$ is a power of $2$, and is divisible by
$4$ otherwise. We prove this. We also prove that the equation has exactly one normalized solution with integer coefficients;
this is not immediate, because computing the coefficients from the equation involves a division by $2$ at every step. The proof
works with the inverse series $R$, which has the form $R=x-2xL$ for a series $L$ that agrees with the generating series
of the Catalan numbers modulo $2$. Modulo $4$, the series $R$ undoes its own substitution, so $A$ agrees with $R$, and the
pattern comes from the parity of the Catalan numbers. The results have been checked by computer and proved formally in
the Lean proof assistant.
\end{abstract}
\begin{center}\begin{minipage}{0.88\linewidth}\small
\textbf{How this paper was produced.} This paper reports the result of an AI research effort directed by
Dev~D.~Goyal. ChatGPT (OpenAI), using the GPT-6 model with Pro reasoning and working autonomously from prompts by
Dev~D.~Goyal, found the proof, wrote its Lean~4 formalization, and wrote a first draft of the paper.
Dev~D.~Goyal chose the problems and asked for proofs, formal verification and revisions, but did not supply proof
ideas. Claude Opus~5.5, used with Claude Code, then checked the work independently: it rebuilt the formal proofs from
source and confirmed that they rely only on Lean's three standard axioms, recomputed the sequence from its definition
compared it with the published terms and tested every claim numerically, and read every step of the proofs. It also wrote the present
version of this paper.
Dev~D.~Goyal has verified the results and is responsible for them. The paper has not been peer reviewed.
The section ``About this paper'' at the end gives the details.
\end{minipage}\end{center}

\section{Introduction}

\subsection{The sequence}\label{sec:sequence}
The \emph{On-Line Encyclopedia of Integer Sequences} (OEIS)~\cite{oeisdb} is a large public database of integer
sequences. Many
entries also record \emph{conjectures}: patterns that the contributor observed in the data but did not prove.

This paper is about entry A389542~\cite{oeis}, contributed by Paul~D.~Hanna in November 2025. It defines a power
series $A(x)=a(1)\,x+a(2)\,x^2+a(3)\,x^3+\cdots$ with integer coefficients and $a(1)=1$ by the equation
\begin{equation}\label{eq:oeis}
A\bigl(A(x)^2-x^2\bigr)=4A(x)^3 .
\end{equation}
The left side is the series obtained by substituting $A(x)^2-x^2=4x^3+\cdots$ for $x$ in $A$
(Section~\ref{sec:background} explains substitution). On the right, $A(x)^3$ is an ordinary power. Theorem~\ref{thm:exist} shows that there is exactly one such
series. Its first coefficients, starting at $n=1$, are 1, 2, 10, 64, 446, 3328, 26004, 209856, 1735734, 14636032, \dots{} The entry reads,
in part:
\begin{quote}\small\ttfamily\raggedright G.f. A(x) satisfies A( A(x)\^{}2 - x\^{}2 ) = 4*A(x)\^{}3.\par\smallskip Conjecture: for n > 1, a(n) == 2 (mod 4) if n = 2\^{}k + 1 (for some k >= 0) otherwise a(n) is divisible by 4.\par\smallskip (3) R(x)\^{}2 = x\^{}2 - R(4*x\^{}3) where R(A(x)) = x.\par\smallskip\hfill\textrm{--- OEIS A389542, Paul D. Hanna, November 2025}\end{quote}

Here ``G.f.''\ stands for \emph{generating function}, the series whose coefficients are the terms of the sequence. To
avoid confusion with the equation numbers of this paper, we refer to the formula numbered (3) in the entry as (E3). In
formula (E3), $R$ is the compositional inverse of $A$: the series that undoes the substitution of $A$
(Section~\ref{sec:background}).

\subsection{The question}
The conjecture concerns the remainders of $a(n)$ on division by $4$. It says that, from $n=2$ on, $a(n)$ is divisible
by $4$, except when $n$ is one more than a power of $2$, that is, when $n$ is one of $2,3,5,9,17,33,\dots$; then $a(n)$
leaves remainder $2$. The table shows the remainders of the first $48$ terms; the exceptional indices are printed in bold.
\begin{center}\small\setlength{\tabcolsep}{3.2pt}
\begin{tabular}{>{\raggedright\arraybackslash}p{6.4em}>{\centering\arraybackslash}p{1.15em}>{\centering\arraybackslash}p{1.15em}>{\centering\arraybackslash}p{1.15em}>{\centering\arraybackslash}p{1.15em}>{\centering\arraybackslash}p{1.15em}>{\centering\arraybackslash}p{1.15em}>{\centering\arraybackslash}p{1.15em}>{\centering\arraybackslash}p{1.15em}>{\centering\arraybackslash}p{1.15em}>{\centering\arraybackslash}p{1.15em}>{\centering\arraybackslash}p{1.15em}>{\centering\arraybackslash}p{1.15em}>{\centering\arraybackslash}p{1.15em}>{\centering\arraybackslash}p{1.15em}>{\centering\arraybackslash}p{1.15em}>{\centering\arraybackslash}p{1.15em}}
\toprule
$n$ & 1 & \hl{2} & \hl{3} & 4 & \hl{5} & 6 & 7 & 8 & \hl{9} & 10 & 11 & 12 & 13 & 14 & 15 & 16 \\ \midrule
$a(n)\bmod 4$ & 1 & \hl{2} & \hl{2} & 0 & \hl{2} & 0 & 0 & 0 & \hl{2} & 0 & 0 & 0 & 0 & 0 & 0 & 0 \\
\bottomrule
\end{tabular}
\\[4pt]
\begin{tabular}{>{\raggedright\arraybackslash}p{6.4em}>{\centering\arraybackslash}p{1.15em}>{\centering\arraybackslash}p{1.15em}>{\centering\arraybackslash}p{1.15em}>{\centering\arraybackslash}p{1.15em}>{\centering\arraybackslash}p{1.15em}>{\centering\arraybackslash}p{1.15em}>{\centering\arraybackslash}p{1.15em}>{\centering\arraybackslash}p{1.15em}>{\centering\arraybackslash}p{1.15em}>{\centering\arraybackslash}p{1.15em}>{\centering\arraybackslash}p{1.15em}>{\centering\arraybackslash}p{1.15em}>{\centering\arraybackslash}p{1.15em}>{\centering\arraybackslash}p{1.15em}>{\centering\arraybackslash}p{1.15em}>{\centering\arraybackslash}p{1.15em}}
\toprule
$n$ & \hl{17} & 18 & 19 & 20 & 21 & 22 & 23 & 24 & 25 & 26 & 27 & 28 & 29 & 30 & 31 & 32 \\ \midrule
$a(n)\bmod 4$ & \hl{2} & 0 & 0 & 0 & 0 & 0 & 0 & 0 & 0 & 0 & 0 & 0 & 0 & 0 & 0 & 0 \\
\bottomrule
\end{tabular}
\\[4pt]
\begin{tabular}{>{\raggedright\arraybackslash}p{6.4em}>{\centering\arraybackslash}p{1.15em}>{\centering\arraybackslash}p{1.15em}>{\centering\arraybackslash}p{1.15em}>{\centering\arraybackslash}p{1.15em}>{\centering\arraybackslash}p{1.15em}>{\centering\arraybackslash}p{1.15em}>{\centering\arraybackslash}p{1.15em}>{\centering\arraybackslash}p{1.15em}>{\centering\arraybackslash}p{1.15em}>{\centering\arraybackslash}p{1.15em}>{\centering\arraybackslash}p{1.15em}>{\centering\arraybackslash}p{1.15em}>{\centering\arraybackslash}p{1.15em}>{\centering\arraybackslash}p{1.15em}>{\centering\arraybackslash}p{1.15em}>{\centering\arraybackslash}p{1.15em}}
\toprule
$n$ & \hl{33} & 34 & 35 & 36 & 37 & 38 & 39 & 40 & 41 & 42 & 43 & 44 & 45 & 46 & 47 & 48 \\ \midrule
$a(n)\bmod 4$ & \hl{2} & 0 & 0 & 0 & 0 & 0 & 0 & 0 & 0 & 0 & 0 & 0 & 0 & 0 & 0 & 0 \\
\bottomrule
\end{tabular}
\end{center}

\subsection{What this paper proves}
\textbf{Main result} (Theorem~\ref{thm:main}). \emph{The conjecture is true: for every $n>1$, $a(n)\equiv2\pmod4$ if
$n=2^k+1$ for some $k\ge0$, and $a(n)\equiv0\pmod 4$ otherwise.}

We also prove that \eqref{eq:oeis} has exactly one normalized solution with integer coefficients
(Theorem~\ref{thm:exist}). This needs an argument: when the coefficients are computed from \eqref{eq:oeis} one at a time,
each new coefficient $a(n)$ first appears in degree $n+1$ with the factor $2$, so each step involves a division by $2$.

The key steps are as follows. The inverse series $R$ satisfies $R(x)^2+R(4x^3)=x^2$ (formula (E3) of the entry). We write
it as $R=x-2xL$; the equation for $R$ then becomes an equation for $L$ that can be solved without any division, and
modulo $2$ it is the equation of the Catalan numbers. Modulo $4$, a short computation shows that $R(R(x))=x$: the
series $R$ undoes its own substitution. Since $A$ also undoes the substitution of $R$, we get $A\equiv R=x-2xL\pmod4$,
so $a(n)\equiv-2\,\coeff{n-1}{L}$, and the parity of the Catalan numbers gives the pattern.

\subsection{What was known}
The entry records the equation satisfied by the inverse series (formula (E3)) and the conjecture. We are not aware of an
earlier proof of the conjecture, or of a proof that the coefficients are integers. The parity of the Catalan numbers,
which we use, is classical~\cite{deutschsagan}.

\subsection{Motivation}
Equations in which an unknown series is substituted into an expression built from itself are a natural source of
integer sequences, and Hanna has contributed many of them to the OEIS. Two questions come up again and again for such
equations: are the solutions integers, and do their remainders follow a pattern? Here both answers rest on the same
observation, that $R$ has the special form $x-2xL$. A series that undoes its own substitution is called an
\emph{involution}; the proof shows that, modulo $4$, the inverse series is one. The final answer depends only
on the binary digits of $n-1$, a kind of pattern that is typical of sequences defined by polynomial equations modulo
$2$~\cite{allouche,christol}.

\subsection{How to read this paper}
Section~\ref{sec:background} explains the notions we need: power series, substitution, compositional inverses, and
remainders. Section~\ref{sec:results} states the results, and Section~\ref{sec:proofs} proves them.
Section~\ref{sec:checks} describes how the results were checked, by computer and by a formal proof in Lean.
Section~\ref{sec:discussion} discusses the scope of the results and open questions.
\section{Background}\label{sec:background}
This section collects the definitions and facts used later. Readers who know formal power series and congruences can skip it and return when needed.
\subsection{Formal power series}
A \emph{formal power series} is an expression
\[
U(x)=u_0+u_1x+u_2x^2+u_3x^3+\cdots ,
\]
determined by its sequence of coefficients $u_0,u_1,u_2,\dots$; see~\cite{niven,stanley}. No number is ever
substituted for $x$, so questions of convergence do not arise. We write $\coeff{n}{U}=u_n$ for the coefficient of $x^n$ in $U$; two series are equal
when all their coefficients are equal.

Series are added coefficient by coefficient and multiplied like polynomials:
\[
\coeff{n}{(UV)}=u_0v_n+u_1v_{n-1}+\cdots+u_nv_0 .
\]
Each coefficient of a product is a finite sum. For example, $(1-x)(1+x+x^2+x^3+\cdots)=1$, so $1+x+x^2+\cdots$ is
the \emph{multiplicative inverse} of $1-x$, written $\frac{1}{1-x}$. More generally, every series $D$ with constant term
$1$ has a unique multiplicative inverse $\frac1D$, whose coefficients are determined one at a time by $D\cdot\frac1D=1$
using only addition, subtraction and multiplication. Whenever a quotient of power series appears in this paper, its
denominator has constant term $1$, so the quotient is again a power series.

A \emph{commutative ring} is a set with an addition, a subtraction and a multiplication that satisfy the usual rules
of arithmetic; examples are $\Z$, $\Q$ and the integers modulo $m$. For a commutative ring $R$ we write $\ps{R}$ for
the set of power series with coefficients in $R$; for instance, writing $\F_2$ for the integers modulo $2$, $\ps{\F_2}$
consists of the series with coefficients modulo $2$. We write
$U(0)$ for the constant term $u_0$ of $U$, and $U=O(x^d)$ when the coefficients of $1,x,\dots,x^{d-1}$ in $U$ are all
zero. Thus $U-V=O(x^d)$ means that $U$ and $V$ have the same coefficient of $x^n$ for every $n<d$. A series with
constant term $0$ and coefficient of $x$ equal to $1$, that is, a series of the form $x+u_2x^2+u_3x^3+\cdots$, is
called \emph{normalized}.

\subsection{Substituting one series into another}
If $V(x)$ has constant term $0$, we can substitute $V$ for $x$ in $U$:
\[
U(V(x))=u_0+u_1V(x)+u_2V(x)^2+u_3V(x)^3+\cdots .
\]
This is called \emph{composition}, also written $U\circ V$. It makes sense because $V(x)^k$ has no terms below $x^k$,
so only the finitely many terms with $k\le n$ affect the coefficient of $x^n$. For example, if $U=x+x^2$ and $V=x^2$, then
$U(V(x))=x^2+x^4$. Composition is associative: $U\circ(V\circ W)=(U\circ V)\circ W$.

Substituting a series into itself repeatedly gives its \emph{iterates}: $U^{\circ2}(x)=U(U(x))$,
$U^{\circ3}(x)=U(U(U(x)))$, and so on, with $U^{\circ0}(x)=x$ and $U^{\circ1}=U$. They must not be confused with
\emph{powers} $U^2=U\cdot U$, $U^3=U\cdot U\cdot U$. (The OEIS sometimes writes $A^k(x)$ for the $k$-th iterate. In
this paper the small circle always means an iterate, and a plain exponent always means a power.)

\subsection{Undoing a substitution}
A normalized series can always be ``undone'' by another normalized series.

\begin{lemma}[Compositional inverse]\label{lem:inverse}
For every normalized series $U$ there is exactly one normalized series $V$ with $U(V(x))=x$, and this $V$ also
satisfies $V(U(x))=x$. If $U$ has integer coefficients, so does $V$.
\end{lemma}

We call $V$ the \emph{compositional inverse} of $U$. The OEIS calls it the \emph{series reversion} of $U$. For example,
the compositional inverse of $x-x^2$ is $C(x)=x+x^2+2x^3+5x^4+14x^5+\cdots$, whose coefficients are the Catalan numbers.
Indeed $C-C^2=x$, which says exactly that substituting $C$ into $x-x^2$ gives $x$.

\begin{proof}
Write $V=x+v_2x^2+v_3x^3+\cdots$ with unknown coefficients. In
$U(V)=V+u_2V^2+u_3V^3+\cdots$, the coefficient of $x^n$ (for $n\ge2$) equals $v_n$ plus terms that involve only
$v_2,\dots,v_{n-1}$. Setting it equal to $0$ therefore determines $v_n$ from the earlier coefficients, using only
additions, subtractions and multiplications. So there is exactly one $V$ with $U\circ V=x$, and it has integer
coefficients if $U$ does. Applying the same construction to $V$ gives a normalized $W$ with $V\circ W=x$. Then
$U=U\circ(V\circ W)=(U\circ V)\circ W=W$, so $V\circ U=V\circ W=x$ as well.
\end{proof}

\subsection{Remainders and congruences}
For an integer $m\ge2$ we write $a\equiv b\pmod m$ (``$a$ is congruent to $b$ modulo $m$'') when $a-b$ is divisible by
$m$, that is, when $a$ and $b$ leave the same remainder after division by $m$. The remainder of $a$ on division by
$m$ is called its \emph{residue} modulo $m$; it is one of $0,1,\dots,m-1$.

Addition and multiplication are compatible with congruences, so the residues $0,1,\dots,m-1$, added and multiplied
modulo $m$, form a ring, written $\Z/m\Z$. When $m=p$ is a prime we also write $\F_p$. In $\F_2=\{0,1\}$ we have $1+1=0$,
so addition and subtraction coincide.

For power series with integer coefficients, $U\equiv V\pmod m$ means that \emph{every} coefficient of $U-V$ is
divisible by $m$. Equivalently, the two series become equal when every coefficient is replaced by its residue; we call
this \emph{reducing} the series modulo $m$. The key fact is the following.

\emph{Reduction modulo $m$ is compatible with all the operations above.} To reduce a sum, a product, a substitution, or
an inverse, we may reduce the ingredients first and then compute in $\Z/m\Z$. This holds because every
coefficient of the result is obtained from finitely many coefficients of the ingredients by adding, subtracting and
multiplying, and these operations are compatible with remainders.

\para{Rings.} Some statements below hold over every commutative ring, for example $\Z$, $\Q$ and $\Z/m\Z$. No
division is used in them; this matters because in $\Z/m\Z$ division is not always possible (for instance, $2$ has no inverse in
$\Z/4\Z$).

\subsection{The Frobenius identity}
\begin{lemma}[Frobenius identity]\label{lem:frobenius}
Let $p$ be a prime. For every power series $U$ with coefficients in $\F_p$, we have $U(x)^p=U(x^p)$.
\end{lemma}

\begin{proof}
In $\F_p$, the binomial theorem simplifies to $(a+b)^p=a^p+b^p$, because every binomial coefficient
$\binom pi$ with $0<i<p$ is divisible by $p$. Also $c^p=c$ for every $c\in\F_p$ (Fermat's little theorem). So
$\bigl(\sum_nu_nx^n\bigr)^p=\sum_nu_n^px^{pn}=\sum_nu_nx^{pn}=U(x^p)$. The first equality holds for polynomials by
repeated use of $(a+b)^p=a^p+b^p$; it holds for power series because each coefficient of $U(x)^p$ involves only finitely
many coefficients of $U$.
\end{proof}

For $p=2$ this says that squaring a series modulo $2$ doubles every exponent: for instance,
$(x+x^3)^2\equiv x^2+x^6\pmod2$. In particular, the square of any series has, modulo $2$, no terms of odd degree.

\subsection{The Catalan numbers modulo 2}
\begin{lemma}\label{lem:catalan}
In every commutative ring the equation $U=x+U^2$ has exactly one solution with $U(0)=0$. Over the integers it is the
Catalan series $C(x)=x+x^2+2x^3+5x^4+14x^5+\cdots$, and
\[
C(x)\equiv x+x^2+x^4+x^8+\cdots=\sum_{k\ge0}x^{2^k}\pmod 2 .
\]
So $\coeff{n}{C}$ is odd exactly when $n$ is a power of $2$.
\end{lemma}

\begin{proof}
The coefficient of $x^n$ in $U^2$ involves only $u_1,\dots,u_{n-1}$, so the equation determines $u_n$ one at a time:
$u_1=1$ and $u_n=\sum_{i=1}^{n-1}u_iu_{n-i}$. Over $\F_2$, the series $S=\sum_kx^{2^k}$ satisfies
$x+S^2=x+\sum_kx^{2^{k+1}}=S$ by Lemma~\ref{lem:frobenius}, so by uniqueness $C\equiv S\pmod2$.
\end{proof}
The parity of the Catalan numbers is classical; see for example~\cite{deutschsagan}. (Here the Catalan number
$\frac1{n}\binom{2n-2}{n-1}$ appears as the coefficient of $x^n$, so ``$n$ a power of $2$'' corresponds to the familiar
statement that the $k$-th Catalan number is odd exactly when $k+1$ is a power of $2$.)

\subsection{Doubling modulo 4}
\begin{lemma}\label{lem:double}
Let $U,V$ be series with coefficients in $\Z/4\Z$. Then $2U=2V$ if and only if $U$ and $V$ agree modulo $2$.
Consequently, if $P$ and $Q$ have constant term $0$ and $P\equiv Q\pmod 2$, then $2H(P)=2H(Q)$ in $\Z/4\Z$ for every
series $H$.
\end{lemma}

\begin{proof}
For residues $u,v$ modulo $4$, $2u=2v$ means that $4$ divides $2(u-v)$, that is, $2$ divides $u-v$. Apply this to
every coefficient. For the second claim, $H(P)$ and $H(Q)$ agree modulo $2$ because reduction is compatible with
substitution.
\end{proof}
In other words, \emph{twice a series, modulo $4$, depends only on the series modulo $2$.} This lets us use information modulo
$2$ to get information modulo $4$, without ever dividing by $2$ in $\Z/4\Z$, which is impossible.

\section{Results}\label{sec:results}

The results involve an auxiliary series $L$, defined in the following lemma (proved in Section~\ref{sec:proofs}).

\begin{lemma}\label{lem:L}
In every commutative ring there is exactly one normalized series $L$ with
\[
L(x)=x+L(x)^2-2x\,L(4x^3).
\]
The solution with integer coefficients satisfies $L\equiv x+L^2\pmod 2$, so $L\equiv\sum_{k\ge0}x^{2^k}\pmod 2$:
the coefficient $\coeff{n}{L}$ is odd exactly when $n$ is a power of $2$.
\end{lemma}

\begin{theorem}[Exactly one integer solution]\label{thm:exist}
Equation \eqref{eq:oeis} has exactly one normalized solution $A$ with integer coefficients. Its compositional inverse is
$R=x-2x\,L$, where $L$ is the integer series of Lemma~\ref{lem:L}, and $R$ satisfies $R(x)^2+R(4x^3)=x^2$.
\end{theorem}

\begin{theorem}[Hanna's conjecture]\label{thm:main}
Modulo $4$, $A\equiv R=x-2x\,L$. Consequently, for every $n>1$,
\[
a(n)\equiv\begin{cases}2\pmod4,&\text{if }n=2^k+1\text{ for some }k\ge0,\\[2pt]0\pmod4,&\text{otherwise.}\end{cases}
\]
\end{theorem}

\section{Proofs}\label{sec:proofs}

\para{The idea in brief.} The integer series $L$ can be constructed directly, because its equation determines it
without division. A direct computation shows that $R=x-2xL$ satisfies formula (E3) of the entry, and hence that its inverse solves
\eqref{eq:oeis}. Uniqueness uses that $2c\ne0$ for every nonzero integer $c$. For the congruence, the doubling lemma
(Lemma~\ref{lem:double}) shows that $R(R(x))\equiv x\pmod 4$, so $A$, the inverse of $R$, agrees with $R$ modulo $4$.

\begin{proof}[Proof of Lemma~\ref{lem:L}]
We show that the right side, in degree $n$, involves only coefficients of $L$ below degree $n$; then the equation
determines the coefficients one at a time, using only additions, subtractions and multiplications. Suppose two normalized
series $F$ and $K$ agree below degree $r$, that is, $F-K=O(x^r)$ with $r\ge2$. Then $F^2-K^2=(F-K)(F+K)=O(x^{r+1})$,
because $F+K$ has no constant term, and $x\bigl(F(4x^3)-K(4x^3)\bigr)=O(x^{3r+1})$. So the right side of the equation
changes only in degrees above $r$, as claimed. Modulo $2$ the term $2x\,L(4x^3)$ vanishes, leaving $L=x+L^2$; by
Lemma~\ref{lem:catalan}, $L\equiv\sum_kx^{2^k}\pmod 2$.
\end{proof}

\begin{proof}[Proof of Theorem~\ref{thm:exist}]
\emph{The equation for $R$.} Using the equation of Lemma~\ref{lem:L} in the last step,
\begin{align*}
R(x)^2+R(4x^3)&=\bigl(x^2-4x^2L+4x^2L^2\bigr)+\bigl(4x^3-8x^3L(4x^3)\bigr)\\
&=x^2+4x^2\bigl(x+L^2-2x\,L(4x^3)-L\bigr)=x^2 .
\end{align*}

\emph{Existence.} Let $A$ be the compositional inverse of $R$ (Lemma~\ref{lem:inverse}); it has integer coefficients.
Substitute $A(x)$ for $x$ in $R(x)^2+R(4x^3)=x^2$. Since $R(A(x))=x$, this gives $x^2+R\bigl(4A(x)^3\bigr)=A(x)^2$, that is,
$R\bigl(4A(x)^3\bigr)=A(x)^2-x^2$. Substituting both sides into $A$ gives $4A(x)^3=A\bigl(A(x)^2-x^2\bigr)$, which is
\eqref{eq:oeis}.

\emph{Uniqueness.} Let $U$ be any normalized integer solution of \eqref{eq:oeis}, and let $V$ be its inverse.
Substituting both sides of \eqref{eq:oeis} into $V$ gives $U(x)^2-x^2=V\bigl(4U(x)^3\bigr)$, and then substituting $V(x)$ for $x$ gives
$x^2-V(x)^2=V(4x^3)$. So $V$ also satisfies $V(x)^2+V(4x^3)=x^2$. Suppose $V\ne R$, and let $r$ be the first degree in
which they differ, by $c\ne0$; then $r\ge2$. In degree $r+1$, the squares differ by $2c$ (because
$V^2-R^2=(V-R)(V+R)$ and $V+R=2x+\cdots$), while $V(4x^3)-R(4x^3)$ has no terms below degree $3r$. So $2c=0$, which is
impossible for a nonzero integer $c$. Hence $V=R$ and $U=A$.
\end{proof}

The uniqueness argument uses that $2c\ne0$; it would fail over a ring in which $2c=0$ for some $c\neq0$, such as
$\Z/4\Z$. This is why Theorem~\ref{thm:exist} is stated for integer coefficients.

\begin{proof}[Proof of Theorem~\ref{thm:main}]
Compute modulo $4$. Since $R=x-2xL\equiv x\pmod 2$, Lemma~\ref{lem:double} gives $2L\bigl(R(x)\bigr)\equiv2L(x)$. Hence
\[
R\bigl(R(x)\bigr)=R-2R\,L(R)\equiv R-2RL=(x-2xL)(1-2L)=x-4xL+4xL^2\equiv x\pmod 4 .
\]
So, modulo $4$, $R$ is its own compositional inverse. Since the compositional inverse is unique
(Lemma~\ref{lem:inverse}), and $A$ is the inverse of $R$, we get $A\equiv R=x-2xL\pmod 4$.

For $n>1$ this gives $a(n)\equiv-2\,\coeff{n-1}{L}\pmod 4$. If $\coeff{n-1}{L}$ is odd, then
$-2\,\coeff{n-1}{L}\equiv2\pmod4$; if it is even, then $-2\,\coeff{n-1}{L}\equiv0\pmod 4$. By Lemma~\ref{lem:L},
$\coeff{n-1}{L}$ is odd exactly when $n-1$ is a power of $2$.
\end{proof}

\section{Checking the results}\label{sec:checks}

A mathematical proof can contain mistakes, so the results were checked in two independent ways.

\para{Numerical check.}
We solved \eqref{eq:oeis} coefficient by coefficient with exact rational arithmetic, up to degree $129$. All
coefficients came out as integers, they reproduce all 22 terms printed in the OEIS entry, and the rule of
Theorem~\ref{thm:main} holds for every $n$ with $1<n<130$.
The computation was done with the PARI/GP computer algebra system, using code that was written separately from the
formal proofs (file \texttt{verification/independent\_checks.gp} in the accompanying repository~\cite{repo}).
A computation can only test finitely many cases, so it is a check on the definitions and on the indexing, not a proof.

\para{Formal proof in Lean.}
Lean~\cite{lean} is a \emph{proof assistant}: a program in which mathematical statements and their proofs are written
in a precise formal language, and every step is checked mechanically down to the basic rules of logic. The theorems of Section~\ref{sec:results} are proved in Lean~4 (version 4.19.0) on top of
Mathlib~\cite{mathlib}, Lean's library of formalized mathematics. The statements are collected in the file \texttt{lean/Entries/A389542.lean} of the repository,
in the namespace \texttt{OEIS.A389542}. The full proofs are in \texttt{lean/Families/SquareReversion.lean}.
The table lists where each result of this paper is proved.
\begin{center}\small\renewcommand{\arraystretch}{1.25}
\begin{tabular}{>{\raggedright\arraybackslash}p{0.52\linewidth}>{\raggedright\arraybackslash}p{0.4\linewidth}}
\toprule
Statement in this paper & Lean theorem(s)\\
\midrule
Lemma~\ref{lem:L} (for integer coefficients), and the parity of $L$ & \texttt{auxiliary\_\allowbreak{}exists\_\allowbreak{}unique}, \texttt{auxiliary\_\allowbreak{}parity} \\
$R=x-2xL$, $R^2+R(4x^3)=x^2$, $A$ and $R$ are mutually inverse, $A$ is normalized & \texttt{auxiliary\_\allowbreak{}definition}, \texttt{inverse\_\allowbreak{}equation}, \texttt{inverse\_\allowbreak{}right}, \texttt{inverse\_\allowbreak{}left}, \texttt{normalized} \\
Theorem~\ref{thm:exist}: $A$ solves \eqref{eq:oeis} and is the only normalized integer solution & \texttt{source\_\allowbreak{}equation}, \texttt{source\_\allowbreak{}exists\_\allowbreak{}unique}, \texttt{source\_\allowbreak{}unique} \\
$R(R(x))\equiv x$ and $A\equiv R=x-2xL$ modulo $4$ & \texttt{inverse\_\allowbreak{}involution}, \texttt{inverse\_\allowbreak{}mod\_\allowbreak{}four}, \texttt{series\_\allowbreak{}mod\_\allowbreak{}four} \\
Theorem~\ref{thm:main}, in three equivalent forms & \texttt{residue\_\allowbreak{}mod\_\allowbreak{}four}, \texttt{exceptional\_\allowbreak{}iff}, \texttt{four\_\allowbreak{}dvd\_\allowbreak{}iff} \\
\bottomrule
\end{tabular}
\end{center}
The Lean project was rebuilt from source, and the axioms on which each of these theorems depends were listed with
Lean's \texttt{\#print axioms} command. Every one depends only on the three standard axioms of Lean's logic
(\texttt{propext}, \texttt{Classical.choice}, \texttt{Quot.sound}); no step is left unproved. A formal proof certifies
exactly the statement that Lean checks. Whether that statement says what the OEIS entry means is a separate question,
so the Lean definitions were also compared with the definition in the OEIS entry
(Section~\ref{sec:sequence}).

\section{Discussion}\label{sec:discussion}

\para{Scope.} Theorem~\ref{thm:main} settles the conjecture completely, for every index. It concerns remainders modulo
$4$ only.

\para{Modulo $8$: an open question.} The remainders modulo $8$ also seem to be governed by the binary digits of $n-1$.
For every $n$ with $2\le n<130$, a computer calculation shows that $a(n)\bmod 8$ equals
\begin{itemize}
\item $2$ if $n-1$ is $1$ or $2$;
\item $6$ if $n-1=2^k$ with $k\ge2$;
\item $4$ if $n-1=2^j+2^k$ with $1\le j<k$ (that is, $n-1$ is even and has exactly two binary digits equal to $1$);
\item $0$ otherwise.
\end{itemize}
We have not tried to prove this. The method of this paper suggests a route: write $A=R+4M$ and continue the computation
modulo $8$.

\section*{About this paper}\label{sec:about}
\para{How the work was done.}
This paper is one of 29 papers on conjectures from the OEIS that came out of an AI research effort directed by
Dev~D.~Goyal~\cite{repo}. The roles were as follows.
\begin{itemize}
\item \emph{Dev~D.~Goyal} started the project, chose which problems to pursue, asked for proofs, formal verification and
revisions, and decided what to publish. Dev~D.~Goyal did not supply proof strategies or intermediate arguments.
\item \emph{ChatGPT} (OpenAI), using the GPT-6 model with Pro reasoning, worked autonomously from those prompts. It
searched the OEIS for suitable conjectures, found the proofs, wrote the Lean~4 formalization, and wrote a first draft
of this paper.
\item \emph{Claude Opus~5.5}, used with Claude Code, checked the work independently. It rebuilt the Lean project from source and confirmed that the theorems listed in Section~\ref{sec:checks} rely only on
Lean's standard axioms; compared the Lean definitions with the OEIS definition; recomputed the sequence with separately
written PARI/GP code, compared it with the terms published in the OEIS, and tested every claim of the paper on many
terms; and read every step of the proofs. It then wrote the present version of the paper, expanding the
explanations and, in places, writing out an argument in more detail or by a more elementary route. The theorems of
Section~\ref{sec:results} and the proof strategies are ChatGPT's. The Discussion section, including the short arguments
in it, was written by Claude Opus~5.5; its observations were checked by computer and are not part of the Lean
formalization.
\end{itemize}
Dev~D.~Goyal has verified the results and is responsible for this paper.

\para{Review status.} This paper has not been peer reviewed, and it has not been submitted to a journal or to the OEIS.

\para{Priority.} An OEIS entry that labels a statement as a conjecture shows only that its contributor did not know
a proof; it does not show that no proof existed. During the research only limited searches for earlier work were
made. No claim of priority is made.

\para{Corrections.} Please report any error in the issue tracker of the repository~\cite{repo}, where the Lean sources and the numerical checks are also
available.

\section*{Acknowledgements}
The sequence, the conjecture, and formula (E3) for the inverse series are due to Paul~D.~Hanna~\cite{oeis}.

\begin{thebibliography}{99}\small
\raggedright
\bibitem{oeis} P.~D.~Hanna, entry A389542 (G.f.\ $A(x)$ satisfies $A(A(x)^2-x^2)=4A(x)^3$), \emph{The On-Line Encyclopedia of Integer Sequences}, November 2025. \url{https://oeis.org/A389542}. Text as of October 2026.
\bibitem{deutschsagan} E.~Deutsch and B.~E.~Sagan, Congruences for Catalan and Motzkin numbers and related sequences, \emph{Journal of Number Theory} 117 (2006), 191--215.
\bibitem{allouche} J.-P.~Allouche and J.~Shallit, \emph{Automatic Sequences}, Cambridge University Press, 2003.
\bibitem{christol} G.~Christol, T.~Kamae, M.~Mend\`es France and G.~Rauzy, Suites alg\'ebriques, automates et substitutions, \emph{Bulletin de la Soci\'et\'e Math\'ematique de France} 108 (1980), 401--419.
\bibitem{oeisdb} N.~J.~A.~Sloane and The OEIS Foundation Inc., \emph{The On-Line Encyclopedia of Integer Sequences}. \url{https://oeis.org}.
\bibitem{stanley} R.~P.~Stanley, \emph{Enumerative Combinatorics}, Volume~2, Cambridge University Press, 1999. (Chapter~5 treats formal power series, composition and compositional inverses.)
\bibitem{niven} I.~Niven, Formal power series, \emph{American Mathematical Monthly} 76 (1969), 871--889.
\bibitem{lean} L.~de~Moura and S.~Ullrich, The Lean~4 theorem prover and programming language, in \emph{Automated Deduction -- CADE 28}, Lecture Notes in Computer Science 12699, Springer, 2021, 625--635.
\bibitem{mathlib} The mathlib Community, The Lean mathematical library, in \emph{Proceedings of the 9th ACM SIGPLAN International Conference on Certified Programs and Proofs (CPP 2020)}, 367--381. Version used: commit \texttt{c44e0c8ee63c}.
\bibitem{repo} D.~D.~Goyal, \emph{Proofs and corrections of OEIS conjectures}, repository, 2026. \url{https://github.com/dvgyl/oeis-a396843-proof}.
\end{thebibliography}
\end{document}
